\section{A hierarchy of probability splits and energy comparisons}
\label{sec:hierarchy}

Throughout this section the disorder strength $h>0$ and the interior energy
$E$ are fixed.  All parameters in the cell construction are fixed as well.
Constants may depend on these choices, including the corner-box exponent
chosen below, but not on the outer torus, the upper thinning level, or the
prescribed artificial-probe mask.  We write
\[
 b_l=3^l,\qquad L_l=\frac{b_l-1}{2}-b,
 \qquad q_l=p_l=\exp(-b_l^6).
\]
Only levels above a sufficiently large fixed $l_0$ will be used.  We increase
$l_0$ whenever necessary to apply all the preceding cell estimates.

\subsection{An iid implementation of the rare cell laws}

On the torus of side $b_N$, choose one uniform lattice offset, independently
of the true iid potentials.  For each $l_0\le l\le N$, partition the torus
into squares of side $b_l$ whose centers are congruent to that offset modulo
$b_l$.  These partitions are nested.  In each tile use the translated
geometry and branch law of \Cref{cell:geometry,cell:branch}, with
$L=L_l$.  The word \emph{ancestor} will always refer to strict containment
in this nested family unless the contrary is stated.

We first record the geometric feature that makes conditioning compatible
across levels.  The corner triples belonging to different tiles and levels
are disjoint except when they coincide exactly.  At a coincident triple
the ordering of its three sites, and the two real plane-basis vectors used
for its box coordinates, are the same.  In fact, the missing corner and
its two companions have fixed displacements from the corresponding two
ends of a tile.  Passing to an extreme child preserves those positions:
\[
 b_l+\frac{b_l-1}{2}-b
       =\frac{b_{l+1}-1}{2}-b.
\]
Consequently, if an ancestor triple meets a descendant tile, it is one of
that descendant's own triples.  Also, the set $V\cup B$ of a descendant
is contained in that of its ancestor, and every descendant center belongs
to the ancestor's $V$.  The latter inclusion will ensure that all
descendant reference-potential changes are interior changes when we use
the parent-cell envelope.

Put $K_0=10^6$.  For each distinct corner triple, give its three true
potentials a single box label as follows.  In each product of polynomial
pieces of $\vartheta$ lying in $(-h,h)^3$, use the fixed linear box
coordinates from \Cref{cell:branch}.  Partition the resulting region into
maximal dyadic cubes whose concentric doubles lie in the region.  Their
inverse images are the corner boxes.  All labels are defined from the
true potentials, but we will expose only specified labels.  At exact
corner coincidences we use the same label, rather than successively
refining it.  Null boundaries may be assigned arbitrarily.

Conditional on a box label, the three true variables are uniform on that
box: the original density is constant, and the coordinate change is
linear with constant nonzero Jacobian.  Labels of disjoint triples are
independent.  If a box has width less than $\xi$, its points lie within
$C\xi$ of the boundary of one of the finitely many piece regions.
Indeed, away from those boundary strips a dyadic cube of width comparable
to $\xi$ containing the point has its concentric double inside the region,
so a maximal cube containing it cannot be smaller.  Thus
\begin{equation}
 \PP(\text{width of a given triple's box}<\xi)
       \le C\xi.
 \label{hier:small-box}
\end{equation}
A level-$l$ tile has \emph{valid own labels} when the widths of all four
corner boxes are at least $b_l^{-K_0}$.  The doubled-box requirement also
provides the polynomial clearance from support endpoints and piece
boundaries required in the cell estimates.

Assign to every tile $U$ of level $l$ an independent selector
$\eta_U\in\{0,1\}$ of success probability $q_l$.  These selectors are
independent of all true potentials.  A tile is \emph{eligible} if it is
selected and has no selected strict ancestor.  Eligible tiles are
pairwise disjoint, even when a selected tile later fails its branch test.
An eligible tile is \emph{allowed} if its own labels and those of every
strict ancestor are valid at their respective levels.

For an allowed tile $U$, let $\mathcal L_U$ be the likelihood of its cell
branch law relative to the true product law with its four corner boxes
fixed.  It changes the law on $V$ and the noncorner single-side
coordinates of $B$, leaves the conditional laws on the four corner
triples unchanged, and leaves all remaining sites of the tile unchanged.
Its root factor is the conditional density centered at the reference
$c(v)$ constructed from the true values on $J$.  Apart from the
allowance gate, the likelihood uses only data in $U$ and its own corner
boxes.  In particular it does not use ancestor labels in any other way.
The density bound of \Cref{cell:branch} gives, uniformly over valid boxes,
\[
 \mathcal L_U\le
       \exp\bigl(C L_l^3+C L_l^2\log(2+L_l)\bigr).
\]
After increasing $l_0$, we have $p_l\mathcal L_U\le1$.

Give each tile a further independent uniform variable $t_U$ on $[0,1]$.
Its acceptance indicator is
\begin{equation}
 A_U=\one_{\{U\text{ eligible and allowed}\}}
             \one_{\{t_U\le p_l\mathcal L_U\}}.
 \label{hier:acceptance}
\end{equation}
The selectors, acceptance uniforms, and offset are auxiliary randomness;
the true potentials are sampled first from their unchanged iid law.

\begin{proposition}[Conditional sampling]
\label{hier:sampling}
Fix the offset and all selector marks.  Expose the box identities of the
eligible tiles and of their ancestors.  Conditional on this information,
the unexplored variables in disjoint eligible tiles are independent.
For an allowed eligible tile $U$ of level $l$, its conditional acceptance
probability is $p_l$, and conditional on acceptance its site variables
have exactly the branch law of \Cref{cell:branch} with the indicated
corner conditioning.

These conclusions remain valid for $U$ after fixing status decisions or
other local choices in disjoint eligible tiles, provided the exposed box
identities are retained in the conditioning.  In particular,
\begin{equation}
 \PP(A_U=1\mid\text{offset})\le q_l p_l.
 \label{hier:incidence}
\end{equation}
The unconditional law of the true potentials remains the iid product
law.  The same construction can be made prospectively by omitting all
selectors and ancestor tests above a specified cutoff level.
\end{proposition}

\begin{proof}
The selector marks determine eligibility without using any potential.
Consider an eligible tile $U$.  Any exposed triple that meets $U$ is
either one of its own triples or an ancestor triple coinciding with one
of them.  To see this also for labels exposed on behalf of another
eligible tile, use nesting: that other tile is disjoint from $U$, and an
ancestor of it that meets $U$ either contains $U$ or is disjoint from its
interior.  In the containing case the corner-coincidence property applies.
Thus no exposed label imposes an additional condition on the interior of
$U$ or on its single-side variables.

The conditional true law in $U$ is consequently precisely the reference
law defining $\mathcal L_U$: the four corner triples are uniform in
their fixed boxes, and all other sites have their product laws.  These
conditional laws factor across eligible tiles, since their site sets
are disjoint and the only repeated labels are identical labels on the
same triple.  For an allowed tile,
\[
 \E[\mathcal L_U\mid\text{selectors and exposed labels}]=1.
\]
Integrating the independent uniform $t_U$ proves that its acceptance
probability is $p_l$.  More generally, for every bounded local test
function $F$,
\[
 \E[F\one_{\{A_U=1\}}\mid\text{selectors and exposed labels}]
   =p_l\E[F\mathcal L_U\mid\text{selectors and exposed labels}].
\]
After division by $p_l$, this is exactly the branch expectation.
Local status decisions in disjoint eligible tiles multiply the
conditional product density by factors on those tiles alone, so do not
alter this conclusion.

The event $A_U=1$ requires its own selector to be one.  Conditional on
eligibility and all required valid labels, the additional acceptance
probability is $p_l$; dropping the other restrictions gives
\Cref{hier:incidence}.  Finally, acceptance only attaches auxiliary
indicators to the initially sampled product field, and never resamples
or changes that field.  Omitting the higher-level rules leaves the same
argument on the truncated hierarchy.
\end{proof}

For a stage $n$, $l_0\le n\le N$, place active cell ports at precisely
the accepted roots of levels at least $n$.  Use their branch strengths
and reference potentials.  Every site without an active cell port has
its true real potential.  Because eligible tiles are disjoint, there is
at most one cell port at a site.  Denote the scattering matrix without
artificial probes by $S_{n,N}$.

We may also prescribe at most one artificial probe per site, absent or
of strength $\tau$ or $\tau\vartheta(v_i)$.  Its mask may depend on the
offset and on external randomness independent of the potentials and
branch decisions.  The profile is always evaluated at the true value,
including at an active cell root.  When comparing stages $n<m$, we
assume
\begin{equation}
 0\le\tau\le\exp(-c_*b_m).
 \label{hier:probe-range}
\end{equation}
This is sufficient for the cell estimate at every removed level and at
its immediate parent.  In particular all constants below are uniform
over these masks.  For notational convenience, expectations denoted by
$\E_{\rm off}$ may include any external independent randomness used in
the geometry or deterministic data prescriptions.

\begin{lemma}[The hidden root law]
\label{hier:hidden}
At an accepted root whose cell port is still open, condition on the
selectors, the needed corner-box identities, all branch acceptance
indicators, and the other true potentials.  Do not condition on the
numerical value of the root tile's acceptance uniform.  The root's true
value then has the form
\[
                 v_0=c(v)+\delta,
\]
where $\delta$ has the centered, cutoff Cauchy law of
\Cref{mat:cutoff}, with the scale of that cell port.  The open response,
all other active references and strengths, and any artificial probe at
this root are independent of this hidden draw under the stated
conditioning.
\end{lemma}

\begin{proof}
The root does not belong to any exposed triple.  By
\Cref{hier:sampling}, acceptance tilts its tile's conditional law by
the branch likelihood.  With all other true values fixed, the only
remaining factor is exactly the specified root density centered at
$c(v)$; the reference is a function of the values on $J$.  Acceptance
indicators of other eligible tiles concern disjoint sites.  Indicators
of ineligible descendants are identically zero and give no further
information about the root.

The open system uses $c(v)$ in place of $v_0$.  The support of the hidden
draw lies on the plateau of $\vartheta$, so its artificial probe, if
present, has constant strength there.  References and strengths in
other active cells depend only on their disjoint tiles.  These facts
prove the final assertion.  Restoring true values at other roots does
not change it, since those values also belong to disjoint eligible
tiles.  Conditioning on the numerical acceptance uniform, in contrast,
could truncate the root density and is expressly excluded.
\end{proof}

\subsection{Thinning passes and admissible data}

Fix $n<m\le N$.  A thinning pass acts on a designated union of level-$m$
tiles, chosen from the offset and any independent geometric randomness.
In that region remove all accepted root ports of levels
$n,\ldots,m-1$, restoring their true real values.  Process these levels
in increasing order.  At a given level interpolate by retaining each
eligible accepted root independently with probability $\theta$, with
$\theta$ decreasing from one to zero.  Use fresh independent thinning
marks; never recompute branch choices.  One may subsequently perform a
second pass on the complementary union of level-$m$ tiles.  Each such
pass is covered separately by the estimates below, with its own
starting system.  A \emph{prefix} may end between levels or at an
intermediate value of $\theta$.

Let $f$ be deterministic finite real scalar or Euclidean-vector site
data, possibly depending on the offset.  Its restriction supplies data
on every port set.  Let $g$ be data fixed at the beginning of the pass
such that, for a deterministic $c_g<\infty$,
\begin{equation}
 \abs{g_i-g_j}^2\le c_g^2\abs{f_i-f_j}^2
                    \quad\text{for all sites }i,j.
 \label{hier:data-control}
\end{equation}
We permit $g$ to depend on selectors, branch statuses, and the
start-of-pass response, and on external independent randomness.  The
precise admissibility requirement is that, at each root closure to
which the energy comparison is applied, $g$ is measurable in the
conditioning that preserves that root's hidden law in
\Cref{hier:hidden}.  In particular it must not inspect still hidden
true values needed for future closures, and it must not depend on the
thinning marks of the pass.  This formulation allows response-dependent
adjustments while retaining the conditional closure identity.

For a level-$l$ tile $U$, let $D(U)$ be its immediate parent, and put
\[
 h_D=\sup_{i,j\in D}\abs{f_i-f_j}.
\]
Define the nonnegative charge
\begin{equation}
 \mathcal C_{n,m}(f)
   =\sum_{l=n}^{m-1} q_lp_l\,
        \E_{\rm off}\sum_{U:\,\operatorname{level}(U)=l}
                       h_{D(U)}^2.
 \label{hier:charge-definition}
\end{equation}
The sum includes all tiles of a level, even when the pass acts on only
part of the torus; this only enlarges the bound.

\begin{theorem}[Energy comparison under thinning]
\label{hier:comparison}
There is a deterministic sequence $\zeta_n\downarrow0$, after replacing
it by a decreasing majorant if necessary, such that for all sufficiently
large $n$, every $n<m\le N$, every
allowed artificial-probe mask, and every pass or prefix just described,
\begin{equation}
 \E\mathcal Q_{\rm end}(g)
 \le \E\mathcal Q_{\rm start}(g)
       +C c_g^2\left(
          \zeta_n\E\mathcal Q_{\rm start}(f)
                    +\mathcal C_{n,m}(f)\right).
 \label{hier:comparison-equation}
\end{equation}
The constants and $\zeta_n$ are independent of $m,N$, the pass region,
the mask satisfying \Cref{hier:probe-range}, and the finite dimensions
of the real vector data.
\end{theorem}

We prove the theorem by first establishing the estimate that controls
the interpolation near $\theta=0$.  In a current level-$l$ system, let
$\mathcal A_l$ be the thinnable roots in the pass region that are still
present.  For $a\in\mathcal A_l$, write $U(a)$ for its level-$l$ tile.
Set
\begin{align}
 Z_l(\theta)
  &=\E\sum_{a\in\mathcal A_l}
        \sum_{j:\,\operatorname{site}(j)\notin D(U(a))}
             \abs{S_{ja}}^2\abs{f_j-f_a}^2,
       \label{hier:exterior-mass}\\
 B_l(f)
  &=q_lp_l\,\E_{\rm off}
         \sum_{U:\,\operatorname{level}(U)=l}h_{D(U)}^2.
       \label{hier:level-charge}
\end{align}
Port data in these formulas are evaluated at their physical sites.
Trivially,
\begin{equation}
               Z_l(\theta)\le2\E\mathcal Q_{\rm current}(f).
 \label{hier:uncapped}
\end{equation}

\begin{lemma}[The parent-alternative cap]
\label{hier:cap}
There are constants $C,C_1$ such that, with
$K_l=\exp(C(1+b_{l+1})^{C_1})$, for $0\le\theta\le1$,
\begin{equation}
 Z_l(\theta)
   \le\theta K_l\E\mathcal Q_{\rm current}(f)
                         +C\theta B_l(f).
 \label{hier:cap-equation}
\end{equation}
The estimate holds during either pass and during every level
interpolation in that pass.
\end{lemma}

\begin{proof}
First replace the datum $f_a$ by that at its parent center $c_D$:
\[
 \abs{f_j-f_a}^2
       \le2\abs{f_j-f(c_D)}^2+2h_D^2.
\]
The additional term is bounded using the unit row mass of $S$ and
\Cref{hier:incidence}.  Each accepted level-$l$ root survives with
probability $\theta$ independently, so the total additional expectation
is at most $C\theta B_l(f)$.  It remains to treat the first term, grouped
by parent $D$.

Fix such a parent in the pass region.  The presence of a thinnable
accepted child implies that all selector marks on $D$ and above it are
zero and that all labels on $D$ and above it are valid.  For comparison,
however, leave the selector of $D$ itself free.  Condition on:
\begin{enumerate}[label=(\roman*),leftmargin=*]
 \item absence of selected strict ancestors of $D$;
 \item the identities of the valid labels of $D$ and its strict
       ancestors;
 \item true potentials outside the parent's $V\cup B$;
 \item exterior selector marks, auxiliary uniforms, and thinning
       outcomes determining choices in tiles disjoint from $D$.
\end{enumerate}
Only this conditioned event can contribute to the upper expression;
outside it that expression is zero.  In particular, do not expose finer
labels, statuses, or true potentials inside the parent.

No eligible tile strictly contains $D$ under this conditioning.
Changing $D$'s selector affects only its descendants.  Exterior
eligibility is unchanged, and any shared ancestor validity tests use
the labels already fixed.  If a label on behalf of an exterior tile
meets the parent, it is one of the fixed parent triples.  Hence the
exterior active ports, their strengths and references, and their
already restored true values are fixed across the alternatives below.
The remaining law on the parent is the true product law with precisely
the parent corner-box conditioning.

Fix an exterior source port $j$ whose physical site is outside $D$.
Schur elimination of the fixed exterior produces the same passive
matrix $T$ and source vector $k$ at the parent's $B$ in both
alternatives.  Artificial probes in the exterior are part of these
fixed data.  All possible descendant centers are in the parent's $V$,
and there are no cell ports on its $B$.  Every active descendant
strength is between $\exp(-C L_{l+1})$ and one, and its reference
potential is bounded, because its level does not exceed $l$.

In the alternative contributing child survivors, apply
\Cref{trans:envelope} to all possible internal active subsets.  The
pointwise envelope furnished there depends only on the true potentials
and the fixed exterior data.  The same pointwise envelope applies to every
internal branch choice and thinning outcome, so all internal selector,
label, and acceptance restrictions may be dropped in its upper bound.
Assign a fresh Bernoulli-$\theta$ survival mark to each of the nine level-$l$ children,
including those not accepted.  A child contribution is possible only
if at least one of these marks is one.  This event has probability at
most $9\theta$, independently of the envelope.  Thus, writing
$\mathcal S_D$ for the thinnable survivors in $D$, the conditional
upper bound is
\begin{equation}
 \E\left[\sum_{a\in\mathcal S_D}\abs{S_{ja}}^2
                              \,\middle|\,\text{fixed data}\right]
 \le\theta\exp\bigl(C(1+b_{l+1})^{C_1}\bigr)
                   \E_{B,\rm true}\rho.
 \label{hier:parent-upper}
\end{equation}
Here $\rho$ is the parent's boundary weight used in
\Cref{trans:cell}; the expectation includes its fixed corner boxes.

In the comparison alternative, $D$ itself is selected and accepted.
It is eligible because its strict ancestor selectors are zero, and
allowed because its labels and those above are valid.  By
\Cref{hier:sampling}, this alternative has conditional probability
$q_{l+1}p_{l+1}$ and its conditional law is the parent branch law.
Its root $a_D$ is still active: the pass is presently at level $l$,
and it has not thinned level $l+1$.  If this is the second pass, the
earlier pass did not touch $D$, since its region and the present region
are unions of level-$m$ tiles and $l+1\le m$.  Therefore the alternative
is an actual member of the current ensemble, not a new system outside
that ensemble.  The lower estimate in \Cref{trans:cell} yields
\begin{equation}
 \E\left[\one_{\{a_D\text{ active}\}}\abs{S_{j a_D}}^2
                              \,\middle|\,\text{fixed data}\right]
 \ge c q_{l+1}p_{l+1}
        \exp\bigl(-C(\log L_{l+1})^2\bigr)
                   \E_{B,\rm branch}\rho.
 \label{hier:parent-lower}
\end{equation}
The harmless normalization of a scattering amplitude is absorbed in
$c$.

The two boundary averages in
\Cref{hier:parent-upper,hier:parent-lower} have identical corner-box
laws.  They differ only in the single-side values and their permitted
artificial absorptions.  The pointwise weight comparison accompanying
\Cref{imp:weighted} compares all such bounded settings by a fixed
factor.  Integrating that comparison shows that the two averages of
$\rho$ are comparable, without imposing an exterior norm bound or
assuming independence after a branch tilt.  Inverting the parent
incidence probability costs
\[
 (q_{l+1}p_{l+1})^{-1}=\exp(2b_{l+1}^6),
\]
which can be absorbed in $K_l$.  Consequently, under the specified
conditioning,
\begin{equation}
 \E\sum_{a\in\mathcal S_D}\abs{S_{ja}}^2
       \le\theta K_l\E\left[
           \one_{\{a_D\text{ active}\}}\abs{S_{j a_D}}^2\right].
 \label{hier:parent-transfer}
\end{equation}
Absent exterior source ports contribute zero on both sides.  Their
status was fixed in the conditioning, so the inequality can be summed
over exterior source indices and then integrated.

Multiply by $\abs{f_j-f(c_D)}^2$ and sum over parents.  For an active
parent this datum is exactly $f_{a_D}$.  Each parent appears only once
at this level, and the resulting oriented parent-to-exterior sum is
at most $2\mathcal Q_{\rm current}(f)$.  Combining it with the
oscillation charge proves \Cref{hier:cap-equation}, after enlarging
$K_l$ by a fixed factor.
\end{proof}

\begin{proof}[Proof of \Cref{hier:comparison}]
Fix a level $l$ during the pass.  Conditional on the initial data and
the true potentials, the expectation over its finitely many thinning
marks is a polynomial in $\theta$.  Differentiation gives a sum of
single-root kept-versus-removed differences, with the other roots
retained independently with parameter $\theta$.  Equivalently, for
$\theta>0$, sum over the roots currently retained and divide by
$\theta$.  The data $g$ are frozen throughout this interpolation.
All scattering energies involved are integrable: unitarity and
\Cref{hier:data-control} bound them by finite deterministic pair-data
oscillations and the finite number of ports.

For one retained thinnable root $a$, use the conditioning of
\Cref{hier:hidden} and the admissibility of $g$.  Restoring its true
potential is exactly the closure operation of
\Cref{mat:energy-closure}.  Its expected increase is at most
\[
                         2\abs{\mu_a}\mathcal M_a(g),
\]
where $\mu_a$ is the conditional cutoff coefficient.  Use
\Cref{hier:data-control} to dominate the squared differences by those
of $f$, with factor $c_g^2$.  After this domination, $g$ plays no role
in any further conditioning or cell integration.

For ports $j$ in the parent $D(U(a))$, the uniform bound on
$\abs{\mu_a}$ from \Cref{mat:cutoff} and row mass give a cost at most
$C c_g^2 h_{D(U(a))}^2$ on occurrence of this root.  Its incidence is
at most $\theta q_lp_l$.  After division by $\theta$, the summed local
cost is at most $C c_g^2 B_l(f)$.

For links to ports outside the parent, apply the relative estimate in
\Cref{trans:cell}.  To justify its use, fix selectors, the needed box
identities, the tile's acceptance, exterior true potentials off its
$V\cup B$, and choices in the disjoint eligible tiles, as well as the
independent thinning outcomes.  Its branch law is unchanged by
\Cref{hier:sampling}.  All other active cell ports and their
reference-potential changes are exterior, since the accepted tile has
no eligible descendants.
Accordingly the source and Schur data are fixed in the cell average.
For deterministic $f$ this gives the upper bound
\[
                C c_g^2 e^{-cL_l}\frac{Z_l(\theta)}{\theta}
\]
for the exterior part of the derivative in the direction of removal.
Writing $F_l(\theta)=\E\mathcal Q_{\rm current}(f)$, the complete
upper derivative bound is therefore
\begin{equation}
 C c_g^2\left(B_l(f)+
                    e^{-cL_l}\frac{Z_l(\theta)}{\theta}\right).
 \label{hier:differential}
\end{equation}

Combine \Cref{hier:uncapped,hier:cap-equation}.  Increasing $K_l$ to be
at least two and using $\min(a,b+c)\le\min(a,b)+c$ gives
\[
 \frac{Z_l(\theta)}{\theta}
       \le \min\left(\frac2\theta,K_l\right)F_l(\theta)
                                                +C B_l(f).
\]
The logarithmic interpolation cost is explicit:
\[
 \int_0^1\min\left(\frac2\theta,K_l\right)\dd\theta
                   =2+2\log(K_l/2)
                   \le C(1+\log K_l).
\]
In particular the additive charge is not multiplied by this logarithm.
Define
\[
 a_l=C e^{-cL_l}(1+\log K_l),\qquad
                    \alpha_n=\sum_{l=n}^{\infty}a_l.
\]
The constants in $K_l$ are fixed independently of $l,m,N$.  Since
$\log K_l$ is polynomial in $b_{l+1}=3b_l$ and $L_l$ is proportional
to $b_l$ for large levels, $\alpha_n$ is finite and tends to zero.
Integrating \Cref{hier:differential} over a whole level, or any prefix
of it, costs at most
\[
 c_g^2\left(a_l\sup_{0\le\theta\le1}F_l(\theta)
                                             +C B_l(f)\right).
\]
At $\theta=0$ use continuity of the finite thinning polynomial.

Let $M$ be the supremum of $\E\mathcal Q(f)$ over all prefixes of the
pass, including all intermediate interpolation parameters.  First take
$g=f$.  Summing the preceding estimates up to any prefix gives
\[
 M\le\E\mathcal Q_{\rm start}(f)
                         +\alpha_n M+C\mathcal C_{n,m}(f).
\]
Choose $n$ sufficiently large that $\alpha_n<1/2$.  Then
\begin{equation}
 M\le\frac{\E\mathcal Q_{\rm start}(f)
                        +C\mathcal C_{n,m}(f)}{1-\alpha_n}.
 \label{hier:prefix-supremum}
\end{equation}
Now apply the same integrated estimates to the given $g$ and substitute
\Cref{hier:prefix-supremum}.  The resulting increase is at most
\[
 c_g^2\left(
       \frac{\alpha_n}{1-\alpha_n}\E\mathcal Q_{\rm start}(f)
              +\frac{C}{1-\alpha_n}\mathcal C_{n,m}(f)\right).
\]
This proves the theorem with a decreasing majorant of
$\alpha_n/(1-\alpha_n)$ as $\zeta_n$.  All bounds are uniform over
the specified levels and masks.  For a complementary pass the same
proof starts with that pass's own initial energy; already performed
changes in the other region are simply fixed exterior data in each
local comparison.
\end{proof}

\subsection{The two useful scales of the additive charge}

\begin{lemma}[Charge tails]
\label{hier:charges}
Put $r_n=\sum_{l=n}^{\infty}q_lp_l$, so that $r_n\to0$.
\begin{enumerate}[label=(\roman*),leftmargin=*]
 \item If $f$ has global lattice Lipschitz constant at most $K$, then
 \[
          b_N^{-2}\mathcal C_{n,m}(f)\le C K^2 r_n,
                         \qquad n<m\le N.
 \]
 \item Suppose $f$ vanishes outside a square of side at most $A R$ and
 has Lipschitz constant at most $K/R$.  If $b_m\le R$ and the square
 is interpreted in a coordinate neighborhood of the torus, then
 \[
                     \mathcal C_{n,m}(f)\le C_A K^2 r_n.
 \]
 Both assertions hold for offset-dependent locations and prescriptions,
 and for real Euclidean-vector data.
\end{enumerate}
\end{lemma}

\begin{proof}
There are $b_N^2/b_l^2$ level-$l$ tiles.  The diameter of an immediate
parent is at most $C b_{l+1}$ in the lattice distance.  In the first
case, $h_D^2\le C K^2 b_{l+1}^2$.  The contribution of this level,
divided by area, is at most
\[
 C K^2 q_lp_l\frac{b_{l+1}^2}{b_l^2}
                                      \le C K^2 q_lp_l.
\]
Sum over $l$.

In the second case, parents missing the support have zero oscillation.
At level $l$, at most $C_A(1+R/b_l)^2$ children have a parent meeting
the support.  Also $h_D^2\le C K^2 b_{l+1}^2/R^2$.  Thus the entire
level contributes at most
\[
 C_AK^2q_lp_l(1+R/b_l)^2\frac{b_{l+1}^2}{R^2}
          \le C_AK^2q_lp_l,
\]
because $b_{l+1}\le b_m\le R$.  Summation proves the second claim.
These are deterministic counts conditional on the offset, so averaging
does not change the bounds.
\end{proof}
